
\documentclass[review]{elsarticle}
\usepackage{lineno,hyperref,amsmath,graphicx,booktabs,siunitx,multirow}\modulolinenumbers[5]

\journal{Acta Astronautica}

\begin{document}

\begin{frontmatter}

\title{Thermochemical Performance Analysis of LMP-103S 
Green Propellant Bipropellant Combinations: 
Quantification of ADN Decomposition Diluent Penalties 
via Cantera Equilibrium Modelling}

\author[bologna]{Prashant Gahlot}
\address[bologna]{Department of Industrial Engineering (DIN), 
Alma Mater Studiorum -- University of Bologna, 
Forl\`{i}, Italy}
\ead{prashant.gahlot@studio.unibo.it}

\begin{abstract}
Ammonium dinitramide (ADN)-based green propellants, 
specifically LMP-103S (HPGP), represent a promising 
hydrazine replacement for in-space propulsion. 
However, their catalytic decomposition produces significant 
inert diluents -- nitrogen (N\textsubscript{2}, 37.9\% 
mole fraction) and water vapour (H\textsubscript{2}O, 
45.5\%) -- that reduce effective combustion performance 
when used as bipropellant oxidisers. 
This study presents a systematic thermochemical analysis 
of LMP-103S as oxidiser with three hydrocarbon fuels 
(ethanol, RP-1, liquid methane) using Cantera 3.2.0 
open-source combustion software with the GRI-Mech 3.0 
reaction mechanism. 
Results are compared against a pure O\textsubscript{2} 
baseline to quantify the diluent performance penalty 
across O/F ratios of 2--10 and chamber pressures of 
10--40 bar. 
Key findings indicate average specific impulse penalties 
of 18.2\%, 26.9\%, and 34.6\% for ethanol, RP-1, and 
liquid methane respectively. 
A Monte Carlo uncertainty analysis (N=500) demonstrates 
model robustness with coefficient of variation below 
4.5\% across all performance metrics. 
Mission-level analysis using the Tsiolkovsky rocket equation 
quantifies propellant mass penalties for representative 
deep-space trajectories, providing design guidance for 
green propellant adoption in volume- and mass-constrained 
spacecraft. 
The results identify ethanol at high O/F ratios as the 
optimal fuel pairing for LMP-103S, achieving Isp penalties 
below 1\% at O/F=10 Additionally, a counter-intuitive volumetric performance advantage 
is identified for ethanol/LMP-103S at high O/F ratios, where 
elevated oxidiser density yields density impulse exceeding 
pure O\textsubscript{2} baseline by 6.8\%.
\end{abstract}

\begin{keyword}
Green propellants \sep LMP-103S \sep ADN \sep 
Bipropellant \sep Thermochemical analysis \sep 
Cantera \sep Diluent penalty \sep Deep-space propulsion
\end{keyword}

\end{frontmatter}

\linenumbers

%% ── 1. INTRODUCTION ────────────────────────────────────────
\section{Introduction}

The global push toward sustainable space access has 
accelerated research into hydrazine-replacement green 
propellants. Environmental and regulatory concerns, 
including potential EU REACH restrictions on hydrazine 
handling, have motivated development of ADN-based 
formulations such as LMP-103S (ECAPS HPGP) and 
AF-M315E (ASCENT) \cite{Larsson2011, Bombelli2020}.

LMP-103S has demonstrated successful in-orbit 
performance through NASA's Green Propellant Infusion 
Mission (GPIM, 2019) and ESA's PRISMA mission, 
establishing its operational readiness as a monopropellant 
\cite{GPIM2019}. 
Its extension to bipropellant configurations -- using the 
oxygen-rich decomposition products as an oxidiser source -- 
offers potential density and toxicity advantages over 
conventional N\textsubscript{2}O\textsubscript{4}/MMH 
systems.

However, the catalytic decomposition of ADN in LMP-103S 
produces a gas mixture dominated by inert species: 
approximately 37.9\% N\textsubscript{2} and 45.5\% 
H\textsubscript{2}O by mole, with only 16.6\% active 
O\textsubscript{2} \cite{Larsson2011}. 
These diluents absorb thermal energy in the combustion 
chamber, increase exhaust molecular weight, and reduce 
effective flame temperatures -- mechanisms that have not 
been systematically quantified for bipropellant 
configurations with hydrocarbon fuels.

Prior thermochemical analyses of green propellant systems 
have primarily employed NASA Chemical Equilibrium with 
Applications (CEA) code, which, while accurate for pure 
propellant species, does not facilitate straightforward 
sensitivity analysis or uncertainty propagation 
\cite{Sharon2021}. 
Furthermore, published studies have focused predominantly 
on monopropellant performance characterisation rather 
than bipropellant diluent effects.

This paper addresses the identified gap by: (1) employing 
Cantera 3.2.0 with GRI-Mech 3.0 to compute equilibrium 
combustion properties for LMP-103S bipropellant systems; 
(2) quantifying diluent performance penalties relative 
to a pure O\textsubscript{2} baseline; (3) conducting 
Monte Carlo uncertainty analysis to assess model 
sensitivity; and (4) translating performance differences 
to mission-level propellant mass penalties for 
representative deep-space trajectories.

%% ── 2. METHODOLOGY ─────────────────────────────────────────
\section{Methodology}

\subsection{LMP-103S Decomposition Model}

LMP-103S (ammonium dinitramide in aqueous solution 
with methanol and water) undergoes catalytic decomposition 
prior to the combustion chamber. 
Following Larsson et al. \cite{Larsson2011}, the 
decomposition product mole fractions are:
N\textsubscript{2}: 0.3789, H\textsubscript{2}O: 0.4554, 
O\textsubscript{2}: 0.1657, 
yielding an effective mixture molecular weight of 
24.11 g/mol and an active O\textsubscript{2} mass 
fraction of 22.0\%.

\subsection{Cantera Equilibrium Model}

Bipropellant combustion is modelled as a constant-enthalpy, 
constant-pressure (HP) equilibrium process using Cantera 
3.2.0 \cite{Cantera2023} with the GRI-Mech 3.0 reaction mechanism (53 species, 325 reactions) \cite{GRI30}.
This approach is thermodynamically equivalent to the 
CEA equilibrium calculation, enabling direct methodological 
comparison.

Fuel species in GRI-Mech are represented as:
ethanol (C\textsubscript{2}H\textsubscript{5}OH) 
via C\textsubscript{2}H\textsubscript{4} + H\textsubscript{2};
RP-1 (C\textsubscript{12}H\textsubscript{24}) 
via 6$\times$C\textsubscript{2}H\textsubscript{4};
methane (CH\textsubscript{4}) directly.

\subsection{Performance Parameters}

Specific impulse is calculated via the characteristic 
velocity method:
\begin{equation}
    I_{sp} = \frac{C_f \cdot c^*}{g_0}
\end{equation}
where the Vandenkerckhove thrust coefficient is:
\begin{equation}
    C_f = \Gamma \sqrt{\frac{2\gamma}{\gamma-1}
    \left[1-\left(\frac{p_e}{p_c}\right)^{(\gamma-1)/\gamma}
    \right]}
\end{equation}
and $\Gamma = \sqrt{\gamma}\left(\frac{2}{\gamma+1}
\right)^{(\gamma+1)/[2(\gamma-1)]}$.
Sea-level expansion is assumed ($p_e = 1$ bar).

\subsection{Uncertainty Analysis}

A Monte Carlo uncertainty analysis (N=500 samples) 
propagates input uncertainties in LMP-103S decomposition 
composition ($\pm$5\%), chamber pressure ($\pm$1 bar), 
and O/F ratio ($\pm$0.2) through the equilibrium model 
to quantify 95\% confidence intervals on performance 
predictions.

\subsection{Mission Trade Study}

Propellant mass requirements for representative missions 
are computed via the Tsiolkovsky rocket equation:
\begin{equation}
    m_{prop} = m_{wet}\left[1 - 
    \exp\left(\frac{-\Delta v}{I_{sp} \cdot g_0}\right)
    \right]
\end{equation}
for a 500 kg spacecraft across five mission scenarios 
spanning LEO to deep space ($\Delta v$ = 150--1500 m/s).

%% ── 3. RESULTS ─────────────────────────────────────────────
\section{Results and Discussion}

\subsection{Thermochemical Performance Comparison}

Figure~\ref{fig:performance} presents chamber temperature and specific 
impulse as functions of O/F ratio for all three fuels with LMP-103S 
and pure O\textsubscript{2} oxidisers at $P_c = 25$ bar. LMP-103S 
consistently yields significantly lower flame temperatures across all 
fuel combinations and O/F ratios examined.

For ethanol fuel, chamber temperatures range from 1128~K to 1965~K 
with LMP-103S, compared to 2631--3655~K with pure O\textsubscript{2} 
— a reduction of 25--69\% depending on O/F ratio. RP-1 exhibits 
similar behaviour (1105--1772~K vs. 3033--3617~K), while methane shows 
the most severe temperature suppression (908--1431~K vs. 2728--3482~K). 
The temperature penalty is most pronounced at low O/F ratios where the 
diluent mass fraction is highest relative to active oxygen content.

Specific impulse values range from 144--188~s for LMP-103S systems 
(methane to ethanol at O/F 2--10) compared to 189--272~s for pure 
O\textsubscript{2} at equivalent conditions, all computed for 
sea-level nozzle expansion ($p_e = 1$~bar).

\begin{figure}[htbp]
  \centering
  \includegraphics[width=\textwidth]{Fig1_performance.png}
  \caption{Thermochemical performance of LMP-103S and pure 
  O\textsubscript{2} oxidisers with three hydrocarbon fuels at 
  $P_c = 25$~bar. Top row: chamber temperature. 
  Bottom row: specific impulse.}
  \label{fig:performance}
\end{figure}

\subsection{Diluent Penalty Quantification}

Figure~\ref{fig:penalty} and Table~\ref{tab:penalty} present the 
systematic diluent penalty analysis. The mean specific impulse penalty 
varies from 18.2\% (ethanol) to 34.6\% (methane), consistent with 
the lower active oxygen content of LMP-103S (16.6~mol\%) versus 
pure O\textsubscript{2} (100~mol\%).

A key finding is the strong O/F dependence of the penalty. For 
ethanol, the $I_{sp}$ penalty decreases monotonically from 39.6\% 
at O/F~=~2 to 0.7\% at O/F~=~10. At oxidiser-rich conditions, 
the increased oxygen availability partially compensates for N\textsubscript{2} 
and H\textsubscript{2}O dilution, approaching pure-O\textsubscript{2} 
performance. This behaviour identifies O/F~=~8--10 as the optimal 
operating regime for ethanol/LMP-103S bipropellant systems.

Methane demonstrates the highest sensitivity to diluent effects 
(mean penalty 34.6\%), attributed to its lower molecular weight 
which amplifies the relative contribution of heavier diluent species 
(N\textsubscript{2}, MW~=~28~g/mol; H\textsubscript{2}O, MW~=~18~g/mol) 
to exhaust molecular weight, reducing the $\sqrt{T/\overline{M}}$ term governing exhaust velocity.

\begin{figure}[htbp]
  \centering
  \includegraphics[width=\textwidth]{Fig2_penalty.png}
  \caption{Performance penalty imposed by ADN decomposition diluents 
  relative to pure O\textsubscript{2} baseline. 
  (a) Combustion temperature reduction. 
  (b) Specific impulse reduction.}
  \label{fig:penalty}
\end{figure}

\begin{table}[htbp]
\centering
\caption{LMP-103S diluent penalty relative to pure O\textsubscript{2} 
baseline at $P_c = 25$~bar.}
\label{tab:penalty}
\begin{tabular}{llrrrrr}
\toprule
Fuel & Metric & O/F=2 & O/F=4 & O/F=6 & O/F=8 & O/F=10 \\
\midrule
\multirow{2}{*}{Ethanol} 
  & $\Delta T_c$ (\%) & 69.1 & 59.6 & 46.8 & 36.0 & 25.3 \\
  & $\Delta I_{sp}$ (\%) & 39.6 & 27.0 & 16.0 & 7.9 & 0.7 \\
\midrule
\multirow{2}{*}{RP-1}    
  & $\Delta T_c$ (\%) & 68.9 & 68.0 & 59.0 & 49.7 & 41.6 \\
  & $\Delta I_{sp}$ (\%) & 42.8 & 36.1 & 26.2 & 18.0 & 11.6 \\
\midrule
\multirow{2}{*}{Methane} 
  & $\Delta T_c$ (\%) & 66.7 & 71.4 & 68.0 & 61.0 & 52.0 \\
  & $\Delta I_{sp}$ (\%) & 45.9 & 42.5 & 36.1 & 28.2 & 20.3 \\
\bottomrule
\end{tabular}
\end{table}

\subsection{Exhaust Composition Analysis}

Figure~\ref{fig:exhaust} presents the equilibrium exhaust composition 
for ethanol fuel with both oxidisers. The LMP-103S panel reveals a 
persistent N\textsubscript{2} fraction of 20--35\% across all O/F 
ratios — inert diluent sourced entirely from ADN decomposition that 
absorbs thermal energy without contributing to thrust generation. 
At O/F~=~6, N\textsubscript{2} constitutes 29.8\% of exhaust mole 
fraction, equivalent to carrying approximately 30\% dead mass through 
the nozzle expansion process. In contrast, the pure O\textsubscript{2} 
panel shows no N\textsubscript{2} and a combustion product distribution 
dominated by H\textsubscript{2}O and CO\textsubscript{2} at high O/F, 
consistent with complete oxidation.

\begin{figure}[htbp]
  \centering
  \includegraphics[width=\textwidth]{fig6a_exhaust_species.png}
  \caption{Equilibrium exhaust mole fraction composition for ethanol 
  fuel with LMP-103S (left) and pure O\textsubscript{2} (right) 
  oxidisers at $P_c = 25$~bar. Note the persistent N\textsubscript{2} 
  fraction in LMP-103S exhaust across all O/F ratios.}
  \label{fig:exhaust}
\end{figure}

\subsection{Effect of Chamber Pressure}

Figure~\ref{fig:pressure} presents the effect of chamber pressure 
on $I_{sp}$ and $c^*$ at fixed O/F~=~6. Increasing chamber pressure 
from 10 to 40~bar improves $I_{sp}$ from 160.0 to 190.4~s for 
ethanol/LMP-103S — a 19.0\% improvement. 

Notably, $c^*$ remains constant at 1240.5~m/s across all pressures 
for LMP-103S systems, confirming that combustion thermochemistry is 
pressure-independent at chemical equilibrium. The $I_{sp}$ improvement 
derives entirely from improved nozzle expansion efficiency at higher 
pressure ratios. This result underscores the importance of maximising 
chamber pressure in LMP-103S bipropellant thruster design.

\begin{figure}[htbp]
  \centering
  \includegraphics[width=\textwidth]{Fig3_pressure.png}
  \caption{Effect of chamber pressure on performance at O/F~=~6. 
  (a) Specific impulse increases with pressure due to improved 
  nozzle expansion. (b) Characteristic velocity $c^*$ is 
  pressure-independent, confirming equilibrium thermochemistry.}
  \label{fig:pressure}
\end{figure}

\subsection{Monte Carlo Uncertainty Analysis}

Figure~\ref{fig:mc} presents Monte Carlo results ($N=500$) for 
the nominal operating condition (ethanol/LMP-103S, O/F~=~6, 
$P_c = 25$~bar). Results yield: $T_c = 1686 \pm 76$~K (CV~=~4.52\%), 
$I_{sp} = 181.7 \pm 3.8$~s (CV~=~2.08\%), $c^* = 1241 \pm 24$~m/s 
(CV~=~1.94\%).

The coefficient of variation below 4.5\% across all performance 
metrics confirms model robustness to realistic input uncertainties. 
The higher CV for chamber temperature (4.52\%) compared to $I_{sp}$ 
(2.08\%) reflects the nonlinear temperature sensitivity to composition 
variation. The 95\% confidence interval for $I_{sp}$ spans 
174.1--189.0~s, providing statistically rigorous performance bounds 
for design-stage analysis.

\begin{figure}[htbp]
  \centering
  \includegraphics[width=\textwidth]{Fig4_montecarlo.png}
  \caption{Monte Carlo uncertainty analysis for ethanol/LMP-103S 
  at O/F~=~6, $P_c = 25$~bar, $N = 500$ samples. Input 
  uncertainties: composition $\pm 5\%$, pressure $\pm 1$~bar, 
  O/F ratio $\pm 0.2$.}
  \label{fig:mc}
\end{figure}

\subsection{Mission-Level Implications}

Figure~\ref{fig:mission} and Table~\ref{tab:mission} present 
propellant mass penalties for a 500~kg spacecraft across 
representative mission scenarios.

Ethanol/LMP-103S incurs minimal penalties: $+0.28$~kg ($+0.7\%$) 
for LEO circularisation to $+1.35$~kg ($+0.5\%$) for GTO-to-GEO 
transfer. These penalties are negligible from a mission design 
perspective and are likely offset by LMP-103S handling and 
storage advantages over conventional oxidisers.

Methane/LMP-103S imposes substantial penalties: $+8.53$~kg 
($+24.4\%$) for LEO and $+40.94$~kg ($+15.9\%$) for 
GTO-to-GEO. For a 500~kg spacecraft, this represents 8.2\% 
of initial wet mass — a significant constraint for 
volume-limited platforms.

\begin{figure}[htbp]
  \centering
  \includegraphics[width=\textwidth]{Fig5_mission.png}
  \caption{Mission-level propellant mass comparison, 
  $m_{wet} = 500$~kg. (a) Extra propellant mass penalty 
  vs mission $\Delta v$. (b) Propellant mass for 
  GTO-to-GEO transfer ($\Delta v = 1500$~m/s).}
  \label{fig:mission}
\end{figure}

\begin{table}[htbp]
\centering
\caption{Propellant mass penalties for LMP-103S vs pure 
O\textsubscript{2}, $m_{wet} = 500$~kg.}
\label{tab:mission}
\begin{tabular}{llrrrr}
\toprule
Mission & $\Delta v$ (m/s) & \multicolumn{3}{c}{Extra propellant mass (kg)} \\
\cmidrule{3-5}
 & & Ethanol & RP-1 & Methane \\
\midrule
LEO Circularisation   & 150  & +0.28 & +4.54  & +8.53  \\
Deep Space Correction & 300  & +0.52 & +8.39  & +15.72 \\
Lunar Orbit Insertion & 900  & +1.12 & +18.30 & +34.02 \\
Mars Orbit Insertion  & 900  & +1.12 & +18.30 & +34.02 \\
GTO to GEO            & 1500 & +1.35 & +22.19 & +40.94 \\
\bottomrule
\end{tabular}
\end{table}

\subsection{Density Impulse Analysis}

Figure~\ref{fig:dimp} presents density impulse 
($I_{sp,d} = I_{sp} \times \rho_{mix}$) as a function of O/F ratio. 
A counter-intuitive finding emerges: at O/F~=~10, ethanol/LMP-103S 
achieves $I_{sp,d} = 221{,}887$~s$\cdot$kg/m\textsuperscript{3}, 
exceeding ethanol/pure O\textsubscript{2} at the same O/F 
(207{,}808~s$\cdot$kg/m\textsuperscript{3}) by 6.8\%.

This reversal arises because LMP-103S mixture density (1240~kg/m\textsuperscript{3}) 
exceeds liquid oxygen density (1140~kg/m\textsuperscript{3}), and 
at high O/F ratios the density advantage outweighs the $I_{sp}$ 
deficit. For volume-constrained spacecraft such as CubeSats where 
tank volume is the primary design constraint, ethanol/LMP-103S at 
high O/F ratios may offer superior volumetric performance to 
conventional O\textsubscript{2}-based systems.

\begin{figure}[htbp]
  \centering
  \includegraphics[width=\textwidth]{fig7_density_impulse.png}
  \caption{Density impulse ($I_{sp} \times \rho_{mix}$) as a 
  function of O/F ratio. Ethanol/LMP-103S exceeds ethanol/pure 
  O\textsubscript{2} at O/F~=~10, demonstrating a volumetric 
  performance advantage for high-density oxidiser applications.}
  \label{fig:dimp}
\end{figure}

\subsection{Model Limitations}
The present analysis employs chemical equilibrium assumption, 
which represents the theoretical maximum performance by assuming 
complete reaction at infinite residence time. Real combustion 
chambers exhibit finite-rate kinetic effects, particularly at 
low chamber pressures and for ADN decomposition intermediates 
not captured by GRI-Mech 3.0. Additionally, fuel species are 
represented by GRI-Mech surrogates (C\textsubscript{2}H\textsubscript{4} 
for ethanol and RP-1) which may introduce small systematic 
deviations in C:H:O ratios. Future work employing ADN-specific 
reaction mechanisms and finite-rate chemistry solvers will 
quantify these deviations.
%% ── 4. CONCLUSIONS ─────────────────────────────────────────
\section{Conclusions}

This study quantifies, for the first time using 
open-source combustion modelling tools, the performance 
penalty imposed by ADN decomposition diluents in 
LMP-103S bipropellant systems. 
Key conclusions are:

\begin{enumerate}
\item ADN decomposition diluents impose mean Isp 
penalties of 18.2\%, 26.9\%, and 34.6\% for ethanol, 
RP-1, and methane fuels respectively, relative to a 
pure O\textsubscript{2} baseline.

\item The penalty is strongly O/F-dependent: 
ethanol/LMP-103S achieves near-parity with pure 
O\textsubscript{2} performance at O/F=10 
(0.7\% Isp penalty), identifying an optimal 
operating regime that minimises diluent effects.

\item Chamber pressure significantly mitigates the 
diluent penalty: increasing from 10 to 40 bar 
improves Isp by approximately 19\% for all fuel 
combinations.

\item Monte Carlo uncertainty analysis confirms 
model robustness with Isp coefficient of variation 
below 2.1\%, validating the equilibrium modelling 
approach for design-stage analysis.

\item At the mission level, ethanol is identified 
as the optimal fuel pairing for LMP-103S bipropellant 
configurations in high-$\Delta v$ applications, 
with propellant mass penalties below 1\% for 
GTO-class missions.
\end{enumerate}

Future work will extend this analysis to finite-rate 
combustion kinetics using ADN-specific reaction 
mechanisms, and integrate the performance model 
into a Digital Twin framework for real-time 
combustion chamber health monitoring.

%% ── REFERENCES ─────────────────────────────────────────────
\bibliographystyle{elsarticle-num}
\begin{thebibliography}{9}

\bibitem{Larsson2011}
A. Larsson, N. Wingborg, 
``Green Propellants Based on Ammonium Dinitramide (ADN),'' 
in \textit{Advances in Spacecraft Technologies}, 
InTech, 2011.

\bibitem{Bombelli2020}
A. Bombelli et al., 
``Green Propulsion Advancement,'' 
\textit{AIAA Propulsion and Energy Forum}, 2020.

\bibitem{GPIM2019}
C.~McLean et al., ``NASA's Green Propellant Infusion Mission,'' 
\textit{AIAA Propulsion and Energy Forum}, 2019.

\bibitem{Cantera2023}
D. G. Goodwin et al., 
``Cantera: An Object-oriented Software Toolkit for 
Chemical Kinetics, Thermodynamics, and Transport Processes,'' 
Version 3.0, 2023. 
doi:10.5281/zenodo.8137090

\bibitem{Sharon2021}
B. M. Sharon Rose, G. Srilochan, 
``Hydrogen Peroxide Based Green Propellants for Future 
Space Propulsion Applications,'' 
M.Sc. Thesis, Politecnico di Milano, 2021.

\bibitem{GRI30}
G. P. Smith et al., 
``GRI-Mech 3.0,'' 
\url{http://www.me.berkeley.edu/gri_mech/}

\end{thebibliography}

\end{document}
